\expected{indiana-code-land-qs}

As described in note~\ref{n:contract-deed}, real estate transactions typically
involve two steps: agreement to a \useterm{contract of sale}, and execution of a
\useterm{deed}. The contract of sale is an ordinary contract, so in theory it
can contain any number of ongoing duties between the buyer and seller. But when
someone sells a house, that person probably wants to be done with the place and
not be tied to it forever. This created a political struggle between real estate
buyers who wanted to place more obligations on sellers in the contract, and
sellers who wanted to be finished with the land once the deed was done.

Eventually, states landed on a political compromise. Under the doctrine of
merger, all obligations under the contract of sale ceased upon execution of the
deed. But parties were allowed to execute deeds that contained six
\term[covenant of title]{covenants of title}, imposing an ongoing duty on the
seller toward the buyer:
\begin{itemize}
\item \textbf{Covenant of seisin}: The seller warrants that the seller owns
the land being sold.
\item \textbf{Covenant of right to convey}: The seller warrants that the seller
has the right to sell the land.
\item \textbf{Covenant against encumbrances}: The seller warrants that no one
else has a property interest in the land, such as a mortgage, other than those
that the seller has disclosed to the buyer.
\item \textbf{Covenant of warranty}: The seller promises to defend the buyer if
someone else lays claim to the land.

\item \textbf{Covenant of quiet enjoyment}: The seller warrants that no one else
will lay claim to the land. (This is generally considered the same as the
covenant of warranty.)\footnote{Be aware that the covenant of quiet enjoyment in
real estate conveyances is not the same as the covenant of quiet enjoyment for
leases, \having{intro-leases-quiet-enjoyment}{discussed in
\mref{intro-leases-quiet-enjoyment}}{discussed in
\mref{intro-leases-quiet-enjoyment}}{which provides different rights}.}

\item \textbf{Covenant of further assurances}: If any documents need to be
signed or actions need to be taken to perfect the buyer's ownership rights, the
seller will take care of them.
\end{itemize}
The covenants of title do not attach to every land transaction, however.
Instead, in most states the parties may choose one of three forms of deed:
\begin{itemize}
\item \term[quitclaim deed]{Quitclaim deed}: The seller makes no warranties,
literally quitting any claim to the land without making further promises.
\item \term[general warranty deed]{General warranty deed}: The seller makes all
six covenants.
\item \term[special warranty deed]{Special warranty deed}: The seller makes all
six covenants, but only with respect to defects in title arising from the
seller's own acts or omissions. In other words, the seller makes no promises
about previous owners' mistakes.
\end{itemize}


